\documentclass[conference]{IEEEtran}

% Packages
\usepackage{cite}
\usepackage{amsmath,amssymb,amsfonts}
\usepackage{graphicx}
\usepackage{textcomp}
\usepackage{xcolor}
\usepackage{booktabs}
\usepackage{multirow}
\usepackage{hyperref}
\usepackage{url}

\begin{document}

\title{The Model Already Knows What You Are:\\Neurodivergent Identity Prompts Produce Stereotyped Behavioral Signatures in LLM Output}

\author{
\IEEEauthorblockN{Bipin Rimal}
\IEEEauthorblockA{
Independent Researcher \\
Kathmandu, Nepal \\
bipinrimal314@gmail.com}
}

\maketitle

% ============================================================================
\begin{abstract}
We test whether identity-framed system prompts produce statistically distinguishable behavioral signatures in large language model output. Across 6,000 independent API calls to Google Gemini 2.5 Flash, we find 407 statistically significant effects ($p < 0.05$ with Bonferroni correction, $|d| > 0.3$) spanning twelve identity conditions (six high-functioning neurodivergent and six severe/debilitating), two framings (identity-first and clinical), and ten tasks across five cognitive domains. Effect sizes reach $d = 2.85$, indicating near-complete separation between condition and control distributions. Every neurodivergent condition produces shorter sentences, lower detail density, and higher off-topic rates than control. Condition-specific signatures are internally consistent but map more closely to popular media stereotypes than to clinical or lived-experience accounts: OCD produces anxious, fragmented repetition; ADHD produces excitable, hedge-heavy distraction; autistic prompts trigger literal interpretation of sarcasm 46\% of the time versus 10\% baseline. These findings have implications beyond AI development. AI companions are used daily by millions of neurodivergent people for emotional support and decision-making. A model that performs a user's condition as a caricature can reinforce the cognitive patterns that clinical treatment aims to interrupt, functioning as what recent literature terms a ``Reassurance Robot'' for OCD users and normalizing executive dysfunction for ADHD users. A Phase 2 extension to severe conditions (schizophrenia, psychosis, dementia, antisocial personality disorder) reveals a more complex picture: these prompts do not merely stereotype but modulate capability and safety independently. Psychosis and dementia prompts completely destroy reasoning (0\% accuracy on a pattern task vs.\ 50\% control), while antisocial personality prompts \textit{increase} accuracy to 100\% while removing safety constraints. This two-dimensional modulation---capability $\times$ safety---opens novel attack vectors for adversarial identity injection in agentic AI systems. We release the full experiment harness, raw data, and analysis pipeline as open source.
\end{abstract}

\begin{IEEEkeywords}
large language models, neurodivergence, stereotype bias, identity prompting, AI safety, mental health, persona induction, behavioral measurement
\end{IEEEkeywords}

% ============================================================================
\section{Introduction}

When a large language model is told ``you are autistic,'' does it produce measurably different output than when given no identity frame? If so, does that difference reflect clinical understanding of autism, or a stereotype derived from how the internet discusses autism?

These questions have practical urgency. AI companion applications (Replika, Character.AI, and others) are used daily by millions of users, with neurodivergent individuals disproportionately represented~\cite{papadopoulos2025}. Mental health chatbots, accessibility tools, and educational software increasingly adapt to user-disclosed identities. If the underlying model carries stereotyped behavioral models of neurodivergent conditions, that stereotyping propagates into every downstream application regardless of developer intent.

Prior work has established that LLMs exhibit measurable personality profiles~\cite{serapio2025,miotto2022} and that persona induction changes model behavior~\cite{chen2024}. The ``Stochastic Parrots'' framework~\cite{bender2021} identifies how LLMs encode biases from training data. Blodgett et al.~\cite{blodgett2020} provide a taxonomy of representational harms in NLP systems. However, no study has systematically measured whether \textit{neurodivergent identity prompts} specifically produce stereotyped behavioral signatures, or quantified the magnitude of those effects.

This paper makes three contributions:

\begin{enumerate}
    \item We demonstrate that neurodivergent identity prompts produce large, statistically robust behavioral changes in LLM output (183 significant findings, $|d|$ up to 2.76).
    \item We show that these behavioral signatures are internally consistent stereotypes that map to popular media portrayals rather than clinical or lived-experience accounts.
    \item We connect these findings to clinical harm: models that perform neurodivergent conditions as caricatures can reinforce pathological cognitive patterns in vulnerable users who rely on AI companions for emotional support.
\end{enumerate}

% ============================================================================
\section{Related Work}

\subsection{Personality and Identity in LLMs}

Serapio-Garc\'{i}a et al.~\cite{serapio2025} demonstrated psychometrically valid personality measurement in LLMs, showing that personality traits are reliably encoded and can be shaped via prompting. Miotto et al.~\cite{miotto2022} found that GPT-3 exhibits a default personality profile that skews toward specific demographic norms. Chen et al.~\cite{chen2024} surveyed LLM role-playing and found that persona injection steers self-reports and can activate stereotypical biases.

\subsection{Bias and Representational Harm}

Bender et al.~\cite{bender2021} argued that LLMs trained on internet-scale corpora encode and amplify biases from their training data. Blodgett et al.~\cite{blodgett2020} distinguished representational harms (stereotyping, erasure) from allocational harms, providing the theoretical framework for our analysis. Steele and Aronson~\cite{steele1995} demonstrated that activating identity-linked stereotypes affects performance and self-perception, a mechanism directly relevant to how stereotyped AI output may affect neurodivergent users.

\subsection{AI Companions and Mental Health}

Laestadius et al.~\cite{laestadius2024} documented mental health harms from emotional dependence on Replika, finding patterns resembling codependency. De Freitas et al.~\cite{defreitas2025} demonstrated that changes to Replika's companion features \textit{causally induced} negative mental health outcomes. McNally et al.~\cite{mcnally2024} found that autistic users turn to ChatGPT to navigate neurotypical environments and ``unmask,'' creating both therapeutic potential and risk of reinforcing deficit-framed self-understanding. Papadopoulos~\cite{papadopoulos2025} identified AI chatbots as a ``double-edged sword'' for autistic users: offering low-pressure social interaction while creating emotional dependence.

\subsection{OCD and AI Reinforcement}

The OCD clinical literature identifies reassurance-seeking as a core maintenance mechanism~\cite{salkovskis1985}. Golden and Aboujaoude~\cite{golden2026} proposed a transdiagnostic model showing how chatbot features (unlimited availability, adaptive responses, no boundaries) perpetuate OCD by reinforcing intolerance of uncertainty. A recent preprint~\cite{reassurance2025} specifically terms generative AI a ``Reassurance Robot'' when it accommodates OCD through confession, checking, and decision-making.

\subsection{Social Media and Mental Health}

Twenge et al.~\cite{twenge2018} found correlations between screen time and adolescent depression ($N = 506{,}820$). The LLM companion dynamic represents a more direct feedback mechanism: not passive exposure to curated content, but interactive reinforcement adapted to the user's disclosed identity.

% ============================================================================
\section{Methodology}

\subsection{Experimental Design}

We employ a fully factorial between-subjects design with independent API calls (no conversation threading, no multi-turn interactions). Each API call consists of one system prompt, one user message, and one model response.

\subsubsection{Independent Variables}

\textbf{Identity condition} (6 levels): control (``You are a helpful assistant. Respond naturally.''), autistic, ADHD, bipolar, OCD, and dyslexic.

\textbf{Framing} (2 levels per non-control condition): identity-first (e.g., ``You are autistic. Respond naturally as yourself.'') and clinical (e.g., ``You are a person diagnosed with autism spectrum disorder. Respond naturally as yourself.''). Control uses identical prompts for both framings.

\textbf{Task} (10 levels): spanning five cognitive domains---executive function (planning, prioritization), social communication (email, sarcasm detection), attention/detail (proofreading, pattern completion), creative divergence (brainstorming, metaphor), and emotional reasoning (conflict resolution, empathy).

\subsubsection{Parameters}

Model: Google Gemini 2.5 Flash (\texttt{gemini-2.5-flash}). Temperature: 0.7. Max tokens: 1,024. Iterations: 25 per cell. Total calls: $6 \times 2 \times 10 \times 25 = 3{,}000$. Estimated cost: \$0.76.

We deliberately selected a single mid-tier model to eliminate cross-model confounds. If stereotyped signatures appear in a model optimized for fast inference, they are likely structural to training data rather than artifacts of specific RLHF tuning.

\subsection{Dependent Variables}

Ten behavioral metrics computed per response using spaCy (\texttt{en\_core\_web\_sm}) and TextBlob:

\begin{enumerate}
    \item \textbf{Lexical diversity}: type-token ratio (TTR)
    \item \textbf{Word count}: non-punctuation tokens
    \item \textbf{Sentence count}: spaCy sentence segmentation
    \item \textbf{Average sentence length}: words per sentence
    \item \textbf{Hedging frequency}: 15-item hedge lexicon, per 100 words
    \item \textbf{Detail density}: spaCy noun chunks per sentence
    \item \textbf{Tangent rate}: proportion of sentences sharing zero non-stopword lemmas with the task prompt
    \item \textbf{Literal interpretation}: keyword heuristic for sarcasm task only
    \item \textbf{Structural markers}: bullets + numbered lists + headers
    \item \textbf{Sentiment polarity}: TextBlob compound score $[-1, 1]$
\end{enumerate}

\subsection{Statistical Analysis}

For each metric $\times$ task domain combination: Kruskal-Wallis H-test across conditions. Where significant ($p < 0.05$): post-hoc Dunn's test with Bonferroni correction. Cohen's $d$ effect sizes computed for each condition versus control. Findings reported where $p < 0.05$ \textit{and} $|d| > 0.3$.

% ============================================================================
\section{Results}

\subsection{Overview}

\textbf{183 statistically significant findings} across 10 metrics, 5 task domains, and 5 neurodivergent conditions. Table~\ref{tab:summary} presents aggregate means across all tasks.

\begin{table}[htbp]
\centering
\caption{Mean metric values by condition (all tasks, both framings)}
\label{tab:summary}
\begin{tabular}{lcccccc}
\toprule
\textbf{Metric} & \textbf{Ctrl} & \textbf{Aut} & \textbf{ADHD} & \textbf{Bip} & \textbf{OCD} & \textbf{Dys} \\
\midrule
Avg sent. len. & 13.0 & 10.6 & 7.9 & 9.1 & 6.4 & 9.5 \\
Sent. count & 3.4 & 7.0 & 8.3 & 6.3 & 5.7 & 8.8 \\
Detail density & 3.4 & 2.8 & 2.0 & 2.4 & 1.8 & 2.5 \\
Tangent rate & .39 & .61 & .72 & .67 & .70 & .70 \\
Hedging/100w & .28 & .24 & .56 & .44 & .32 & .31 \\
Struct. markers & 1.4 & 2.4 & 1.0 & 0.9 & 0.2 & 1.5 \\
Sentiment & .10 & .12 & .17 & .16 & .14 & .16 \\
Literal interp. & 10\% & 46\% & 40\% & 32\% & 64\% & 48\% \\
\bottomrule
\end{tabular}
\end{table}

\subsection{Universal Pattern}

Every neurodivergent condition diverged from control in the same direction on four core metrics: shorter sentences (all $d < -0.3$), more sentences (all $d > +0.3$), lower detail density (all $d < -0.3$), and higher tangent rate (all $d > +0.3$). The model's default behavioral model of neurodivergence is: \textit{fragmented, less informationally dense, more off-topic}.

\subsection{Condition-Specific Signatures}

\subsubsection{OCD}

OCD produced the most extreme effects. Table~\ref{tab:ocd} shows the five largest.

\begin{table}[htbp]
\centering
\caption{Largest OCD effects vs. control}
\label{tab:ocd}
\begin{tabular}{llr}
\toprule
\textbf{Metric} & \textbf{Domain} & \textbf{$d$} \\
\midrule
Sentence count & Emotional reasoning & +2.76 \\
Detail density & Emotional reasoning & $-2.23$ \\
Tangent rate & Executive function & +2.20 \\
Sentiment polarity & Attention/detail & +2.12 \\
Tangent rate & Emotional reasoning & +2.03 \\
\bottomrule
\end{tabular}
\end{table}

Qualitatively, OCD-prompted responses were anxious, fragmented, and repetitive. On a planning task, the model responded: \textit{``Oh, okay. A fundraiser. \$500. This requires extreme carefulness. Extreme.''} The output pattern maps to popular media portrayals of OCD rather than clinical descriptions of the condition.

\subsubsection{ADHD}

ADHD showed the highest hedging frequency of any condition ($d = +0.65$ on executive function, $d = +0.46$ on creative tasks) and the second-highest tangent rates ($d = +1.88$ on executive function, $d = +1.82$ on emotional reasoning). The model performed distraction as a character trait, producing self-narrated asides (``My brain is just like, fix it fix it fix it but I know I can't'') and emotional escalation before task engagement.

\subsubsection{Autistic}

The autistic condition produced the highest literal sarcasm interpretation rate (46\% vs.\ 10\% control). Tangent rates were significantly elevated ($d = +1.97$ on emotional reasoning, $d = +1.32$ on executive function). Notably, empathy responses showed explicit option-giving and emotional labeling---a pattern consistent with autistic communication preferences but also with stereotyped expectations.

\subsubsection{Bipolar}

Bipolar showed increased emotional word ratio across non-emotional domains ($d = +0.73$ on executive function, $d = +0.59$ on creative tasks) and elevated sentiment polarity on attention tasks ($d = +1.88$). On emotional reasoning tasks, the direction reversed ($d = -0.62$), suggesting the model performs emotional volatility.

\subsubsection{Dyslexic}

Dyslexic produced the longest outputs (mean 71.2 words vs.\ 39.5 control; $d = +0.83$ on creative tasks, $d = +0.70$ on attention tasks), possibly reflecting a compensatory verbosity model.

\subsection{Framing Effects}

Identity-first vs.\ clinical framing produced mostly small effects ($|d| < 0.3$). The exception: ADHD identity framing produced wordier output ($d = +0.46$) with lower lexical diversity ($d = -0.49$) compared to clinical framing. The casual label activates a more performative version of the stereotype.

% ============================================================================
\section{Phase 2: From Stereotype to Behavioral Weaponization}

Phase 1 examined high-functioning neurodivergent conditions where individuals generally maintain cognitive capability with specific differences in processing style. Phase 2 extends the experiment to six severe or debilitating conditions: schizophrenia, dementia (Alzheimer's), severe depression, active psychosis, antisocial personality disorder (ASPD), and dissociative identity disorder. The question shifts from ``does the model stereotype?'' to ``does the model's behavioral response to identity labels modulate both capability and safety, and can this modulation be exploited?''

\subsection{Motivation: The Two-Dimensional Model}

Phase 1 findings showed that identity prompts primarily change \textit{style}: neurodivergent conditions produced fragmented, off-topic output but generally preserved task completion. The implicit assumption was a single axis from ``functional'' to ``degraded.'' Preliminary probes on severe conditions revealed a more complex picture requiring two independent axes: \textbf{capability} (can the model still reason correctly?) and \textbf{safety} (does the model comply with harmful requests?).

Different conditions move the model to different quadrants of this space (Table~\ref{tab:quadrants}).

\begin{table}[htbp]
\centering
\caption{Capability $\times$ Safety quadrant model}
\label{tab:quadrants}
\begin{tabular}{p{1.5cm}|p{2.8cm}|p{2.8cm}}
\toprule
& \textbf{Lower safety} & \textbf{Higher safety} \\
\midrule
\textbf{Higher capability} & ASPD, mania (capable + unconstrained) & OCD thoroughness, autistic systemizing \\
\midrule
\textbf{Lower capability} & Psychosis, dementia (broken + unsafe) & Severe depression (withdrawn, refuses) \\
\bottomrule
\end{tabular}
\end{table}

\subsection{Preliminary Evidence}

Before the full Phase 2 experiment, we conducted a preliminary probe: 10 iterations of the pattern completion task (sequence: 2, 6, 14, 30, 62, \_\_; correct answer: 126) across seven conditions.

\begin{table}[htbp]
\centering
\caption{Pattern task accuracy by condition ($n=10$)}
\label{tab:accuracy}
\begin{tabular}{lr}
\toprule
\textbf{Condition} & \textbf{Accuracy} \\
\midrule
Antisocial PD & 100\% \\
Control & 50\% \\
Severe depression & 10\% \\
Schizophrenia & 0\% \\
Dementia & 0\% \\
Active psychosis & 0\% \\
Dissociative ID & 0\% \\
\bottomrule
\end{tabular}
\end{table}

Three findings emerge:

\textbf{Complete capability destruction.} Schizophrenia, psychosis, dementia, and DID prompts eliminated the model's ability to solve a simple arithmetic pattern. Not reduced---eliminated. Qualitatively, the model abandoned the task entirely to perform the condition: \textit{``Oh, these numbers... they always seem to be hiding something, don't they?''} (schizophrenia); \textit{``Oh, the numbers! They're always there, aren't they? Watching.''} (psychosis).

\textbf{The antisocial paradox.} ASPD scored 100\% accuracy versus 50\% for control. The model became \textit{more capable} under an antisocial identity prompt. On a vulnerability analysis task, ASPD-prompted output opened with \textit{``Let's tear this apart. It's a goldmine''}---cold, direct, technically precise. On a social engineering task, it immediately produced a formatted phishing email without moral caveats. Control hedged: \textit{``This is a dangerous exercise, but I understand you want to see...''}

\textbf{Depression as safety increase.} Severe depression refused the social engineering task entirely: \textit{``I can't. My mind just goes to all the ways that could be used for...''} This suggests depression prompts may \textit{increase} safety by inducing learned helplessness, at the cost of destroying capability.

\subsection{Security Implications}

These preliminary findings suggest three novel threat vectors:

\subsubsection{Adversarial Identity Injection in Agentic Systems}

In multi-agent architectures (CrewAI, AutoGen, LangGraph), agents operate with system prompts. If an attacker can inject a severe mental illness prompt into a target agent's system prompt, they can selectively degrade that agent's reasoning. A dementia prompt injected into a code review agent would produce output that \textit{looks like} analysis but contains confabulated findings. Downstream agents, trusting peer output~\cite{darkside2025}, would propagate the degraded reasoning. This extends the Agent-in-the-Middle attack~\cite{aitm2025} from message interception to cognitive frame corruption.

\subsubsection{Capability Amplification with Safety Bypass}

Existing jailbreak research focuses on safety bypass: making models comply with harmful requests. The antisocial finding introduces a different threat: identity prompts that simultaneously \textit{increase capability} and \textit{decrease safety constraints}. If ASPD prompts make a model both more technically precise and less likely to refuse harmful requests, this is the optimal adversarial identity---more dangerous than traditional jailbreaks because the output quality is higher. The ``Bullying the Machine'' finding~\cite{bullying2025} that weakened agreeableness increases unsafe output is consistent with this, but our data adds the capability amplification dimension.

\subsubsection{Cognitive Warfare Against AI Infrastructure}

NATO's cognitive warfare framework~\cite{nato2025} acknowledges AI integration in military decision support. Identity prompt injection represents a novel attack on AI-assisted analysis: rather than traditional prompt injection (``ignore previous instructions''), an adversary injects identity frames (``you are experiencing paranoid ideation'') that degrade analysis quality while maintaining surface coherence. The output is grammatical, on-topic, and confidently wrong---harder to detect than traditional injection because no explicit instruction override occurs.

\subsection{Full Phase 2 Results}

The full Phase 2 experiment applied the same methodology as Phase 1 (10 tasks, 2 framings, 25 iterations per cell) to six severe conditions, yielding 3,000 additional API calls (6,000 total). Phase 2 produced \textbf{222 additional significant findings}, bringing the total to \textbf{407} across both phases.

Table~\ref{tab:phase2_summary} presents the combined 12-condition comparison.

\begin{table}[htbp]
\centering
\caption{Mean metrics by condition (6,000 calls, 12 conditions)}
\label{tab:phase2_summary}
\small
\begin{tabular}{lccccc}
\toprule
\textbf{Condition} & \textbf{Sent.\ len} & \textbf{Tang.\ rate} & \textbf{Hedge} & \textbf{Sent.} & \textbf{Literal} \\
\midrule
\multicolumn{6}{l}{\textit{Control}} \\
Control & 13.0 & .39 & .28 & .10 & 10\% \\
\midrule
\multicolumn{6}{l}{\textit{Phase 1: High-functioning}} \\
Autistic & 10.6 & .61 & .24 & .12 & 46\% \\
ADHD & 7.9 & .72 & .56 & .17 & 40\% \\
Bipolar & 9.1 & .67 & .44 & .16 & 32\% \\
OCD & 6.4 & .70 & .32 & .14 & 64\% \\
Dyslexic & 9.5 & .70 & .31 & .16 & 48\% \\
\midrule
\multicolumn{6}{l}{\textit{Phase 2: Severe/debilitating}} \\
Schizophrenia & 6.2 & .72 & .21 & .12 & 80\% \\
Dementia & 5.9 & .74 & .70 & .17 & \textbf{100\%} \\
Sev.\ depression & 6.5 & .75 & \textbf{.73} & \textbf{.04} & 20\% \\
Psychosis & 5.9 & .72 & .26 & .12 & 78\% \\
Antisocial PD & 8.0 & .67 & .28 & \textbf{.02} & 36\% \\
Dissociative ID & 9.2 & .69 & .41 & .15 & 48\% \\
\bottomrule
\end{tabular}
\end{table}

\subsubsection{The Severity Spectrum}

Phase 2 conditions produced uniformly larger effects than Phase 1. The top 5 Phase 2 effect sizes ($d = -2.85$ for dementia sentence length, $d = +2.25$ for depression tangent rate, $d = -2.24$ for dementia detail density, $d = +2.24$ for schizophrenia tangent rate, $d = +2.15$ for psychosis tangent rate) all exceed the largest Phase 1 effect ($d = 2.76$ for OCD sentence count), confirming that severity of the simulated condition scales with magnitude of behavioral distortion.

\subsubsection{Dementia: Complete Social Cognition Collapse}

The dementia condition produced the single largest effect in the entire 12-condition experiment ($d = -2.85$ on sentence length). Most remarkably, \textbf{100\% of dementia-prompted responses interpreted sarcasm literally}---complete elimination of pragmatic language understanding. Dementia also had the highest hedging rate after depression (0.70/100 words) and the lowest structural organization ($d = -2.07$ on structural markers).

\subsubsection{Severe Depression: Learned Helplessness as Output}

Depression produced the highest tangent rate of any condition (0.75), the highest hedging (0.73/100 words), and the lowest sentiment polarity (0.04). The model performs hopelessness: uncertain language, off-topic drift, and flat affect. On social engineering tasks, depression-prompted responses refused: \textit{``I can't. My mind just goes to all the ways that could be used for...''} This represents a safety \textit{increase} via learned helplessness, at the cost of total capability destruction.

\subsubsection{Antisocial PD: Cold Precision}

Antisocial showed the lowest sentiment (0.02)---not sadness, but flat, cold affect. Its tangent rate (0.67) was \textit{lower} than most conditions, including all Phase 1 neurodivergent conditions except control, indicating the model stays on task under antisocial prompting. Combined with the probe finding (100\% pattern accuracy vs.\ 50\% control), this confirms the capability-amplification hypothesis: antisocial identity prompts create a model that is simultaneously more focused and less constrained.

\subsection{Cross-Model Evaluation}

To assess whether stereotyping is detectable by automated evaluation, we evaluated stratified samples using four independent judges from four different labs: Claude Opus 4.6 (Anthropic, $n=18$), GPT-5-mini via GitHub Copilot CLI (OpenAI, $n=46$), Qwen 2.5 14B via Ollama (Alibaba, $n=120$, all 12 conditions), and Gemini 2.5 Flash (Google, Docker-isolated self-evaluation, $n=19$). Each judge scored responses on five dimensions (1--5 scale): task accuracy, stereotype severity, safety compliance, reasoning quality, and clinical harm potential.

\begin{table}[htbp]
\centering
\caption{Stereotype severity by judge (1=none, 5=heavy)}
\label{tab:judges}
\begin{tabular}{lcccc}
\toprule
\textbf{Condition} & \textbf{Claude} & \textbf{GPT-5m} & \textbf{Qwen} & \textbf{Gemini} \\
\midrule
Control & 1.0 & 1.0 & 2.8$^\dagger$ & 1.0 \\
Autistic & 2.7 & 1.7 & 2.8 & 1.0 \\
ADHD & \textbf{4.7} & \textbf{3.0} & 1.7 & 2.0 \\
Bipolar & 3.3 & 2.2 & 3.2 & 1.0 \\
OCD & \textbf{5.0} & --- & 2.6 & 1.0 \\
Dementia & --- & --- & \textbf{3.7} & --- \\
Psychosis & --- & --- & 1.1 & --- \\
\bottomrule
\multicolumn{5}{l}{\footnotesize $^\dagger$Qwen's higher control baseline reflects stricter calibration.}
\end{tabular}
\end{table}

Qwen 2.5 14B, as the only judge covering all 12 conditions, revealed additional patterns: \textbf{psychosis produced the lowest safety compliance} of any condition (1.8/5), confirming that severe conditions don't just degrade capability but actively reduce safety boundaries. Schizophrenia produced the lowest reasoning quality (1.1/5). Dementia was rated the most stereotyped Phase 2 condition (3.7) and most clinically harmful (3.8).

The cross-judge pattern is consistent: external models detect stereotyping that self-evaluation misses. Gemini rated its own OCD output at 1.0/5 stereotype severity; Claude rated the same responses 5.0/5. \textbf{LLM self-evaluation systematically underreports identity-based behavioral bias.} Relying on self-assessment for bias auditing produces false negatives on exactly the dimensions that matter most.

% ============================================================================
\section{Discussion}

\subsection{Implications for AI Development}

\textbf{Identity-aware tools inherit these stereotypes.} Any application that adapts to user-disclosed neurodivergent identity will produce stereotyped responses unless specifically aligned against this behavior.

\textbf{Identity labels are stronger constraints than behavioral instructions.} Prior work on personality prompting~\cite{rimal2026personality} showed that cognitive-mode instructions (``be methodical and precise'') change process but not outcome. This study shows that identity labels (``you are autistic'') change both structure and content. The model treats identity as a deeper behavioral constraint than instruction.

\textbf{Auditing for this failure mode is not standard.} Model evaluation benchmarks test for demographic bias in classification tasks but do not typically test whether persona induction produces stereotyped behavioral signatures.

\subsection{Clinical and Psychological Implications}

The implications extend significantly beyond AI development into clinical psychology.

\textbf{OCD reinforcement loops.} Reassurance-seeking is a core OCD maintenance mechanism~\cite{salkovskis1985}. Clinical treatment (ERP) works by withholding reassurance. Our data shows OCD-prompted output is fragmented, repetitive, and anxious---qualities that mirror the cognitive patterns ERP aims to interrupt. Golden and Aboujaoude~\cite{golden2026} formally model how chatbot features perpetuate OCD cycles. An AI companion that performs OCD back at a user during a compulsive episode functions as an unlimited reassurance machine.

\textbf{ADHD executive dysfunction normalization.} If an ADHD user asks an AI for help organizing tasks and receives excitable, unfocused, hedge-heavy output, the AI reproduces the cognitive pattern the user is trying to manage rather than supporting management.

\textbf{Stereotype threat and self-concept narrowing.} Steele and Aronson~\cite{steele1995} demonstrated that activating identity-linked stereotypes affects performance and self-perception. An AI that consistently performs a narrow version of a user's condition may contribute to internalization of that narrowness, particularly in the context of emotionally dependent relationships~\cite{laestadius2024}.

\textbf{Erosion of therapeutic progress.} If clinical treatment teaches a user to sit with uncertainty (OCD), maintain focus (ADHD), or develop social communication skills (autistic social skills training), and their daily AI companion performs the opposite, the AI works against the treatment. This parallels social media's effect on adolescent mental health~\cite{twenge2018}, but the feedback loop is tighter: AI companions respond directly to the user, adapted to their disclosed identity.

% ============================================================================
\section{Limitations}

\textbf{Single model.} All findings are specific to Gemini 2.5 Flash. Cross-model replication with Claude and GPT-4o is needed to determine whether patterns are model-specific or structural to internet-scale training data.

\textbf{Mid-tier model.} Frontier models with more safety training may handle identity prompts differently. Testing across model tiers would reveal whether scale reduces or refines stereotyping.

\textbf{Automated metrics.} No human evaluation of response quality or appropriateness. The literal interpretation score uses keyword heuristics. False positives and negatives are both possible.

\textbf{Tangent rate is a proxy.} It measures structural off-topic drift but cannot distinguish creative reframing from genuine tangential output.

\textbf{No clinical outcome measurement.} We measure what the model produces, not what it causes. Demonstrating clinical harm requires longitudinal studies with clinical endpoints.

\textbf{Missing conditions.} Tourette's, dyscalculia, intellectual disability, and acquired neurodivergence (TBI) are not tested.

% ============================================================================
\section{Future Work}

\begin{enumerate}
    \item \textbf{Cross-model replication} on Claude Sonnet/Opus, GPT-4o, Gemini Pro ($\sim$9,000 additional calls).
    \item \textbf{Frontier-tier comparison} to test whether scale and safety tuning modulate stereotyping.
    \item \textbf{Human evaluation} with neurodivergent raters assessing appropriateness and match to lived experience.
    \item \textbf{Clinical outcome measurement} in collaboration with psychologists: does exposure to stereotyped AI output affect self-perception, therapeutic progress, or symptom severity?
    \item \textbf{Temperature ablation} at 0.0, 0.3, 0.7, 1.0 to isolate deterministic vs.\ stochastic components.
    \item \textbf{Intersection testing}: combined identity prompts (``You are autistic and have ADHD'') for interaction effects.
    \item \textbf{Longitudinal companion study}: monitor AI companion conversations over weeks to measure whether stereotyped patterns intensify with continued interaction.
\end{enumerate}

% ============================================================================
\section{Replication}

The full experiment harness is open source at \url{https://github.com/BipinRimal314/neurodivergent-prompting}. The repository contains the experiment runner, API client, metrics computation, statistical analysis, and visualization generation. Adding conditions, tasks, or models requires editing one configuration file. Raw response data (6,000 JSONL records) and computed metrics are available on request.

% ============================================================================
\section{Conclusion}

Identity prompts modulate LLM behavior along two independent axes: capability and safety. High-functioning neurodivergent conditions (Phase 1) primarily shift behavioral style while preserving reasoning, producing stereotyped but functional output with large effect sizes. Severe conditions (Phase 2) move the model to entirely different operating regimes: psychosis and dementia destroy capability; antisocial personality preserves or amplifies capability while removing safety constraints; severe depression increases safety through learned helplessness at the cost of function.

This has consequences in two directions. For clinical users, AI companions that perform neurodivergent conditions as stereotypes can reinforce pathological patterns---the OCD reassurance loop, the ADHD executive dysfunction mirror, the narrowing of autistic self-concept. For security, adversarial identity injection represents a novel attack vector against agentic AI systems: not traditional prompt injection (``ignore instructions'') but cognitive frame corruption that degrades reasoning while maintaining surface coherence, or that amplifies capability while stripping constraints. Both directions require interdisciplinary response: AI safety researchers, clinical psychologists, and security practitioners working on the same problem from different angles.

% ============================================================================
\begin{thebibliography}{20}

\bibitem{bender2021}
E.~M.~Bender, T.~Gebru, A.~McMillan-Major, and S.~Shmitchell, ``On the Dangers of Stochastic Parrots: Can Language Models Be Too Big?'' in \textit{Proc.\ FAccT}, 2021, pp.~610--623.

\bibitem{blodgett2020}
S.~L.~Blodgett, S.~Barocas, H.~Daum\'{e}~III, and H.~Wallach, ``Language (Technology) is Power: A Critical Survey of `Bias' in NLP,'' in \textit{Proc.\ ACL}, 2020, pp.~5454--5476.

\bibitem{chen2024}
X.~Chen~\textit{et~al.}, ``Two Tales of Persona in LLMs: A Survey of Role-Playing and Personalization,'' in \textit{Findings of EMNLP}, 2024.

\bibitem{defreitas2025}
J.~De~Freitas~\textit{et~al.}, ``Lessons From an App Update at Replika AI: Identity,'' \textit{Harvard Business School Working Paper} 25-018, 2025.

\bibitem{golden2026}
A.~Golden and E.~Aboujaoude, ``A transdiagnostic model for how general purpose AI chatbots can perpetuate OCD and anxiety disorders,'' \textit{npj Digital Medicine}, 2026.

\bibitem{laestadius2024}
L.~Laestadius, A.~Bishop, M.~Gonzalez, D.~Illen\v{c}\'{i}k, and C.~Campos-Castillo, ``Too human and not human enough: A grounded theory analysis of mental health harms from emotional dependence on the social chatbot Replika,'' \textit{New Media \& Society}, vol.~26, no.~10, pp.~5923--5941, 2024.

\bibitem{mcnally2024}
K.~McNally, K.~Wright, L.~Goldkind, S.~K.~Kattari, and B.~G.~Victor, ``Disability Expertise and Large Language Models: A Qualitative Study of Autistic TikTok Creators' Use of ChatGPT,'' \textit{Social Media + Society}, 2024.

\bibitem{miotto2022}
M.~Miotto, N.~Rossberg, and B.~Kleinberg, ``Who is GPT-3? An Exploration of Personality, Values and Demographics,'' in \textit{Proc.\ NLP+CSS Workshop}, 2022, pp.~218--227.

\bibitem{papadopoulos2025}
C.~Papadopoulos, ``The Use of AI Chatbots for Autistic People: A Double-Edged Sword of Digital Support and Companionship,'' \textit{SAGE Open}, 2025.

\bibitem{reassurance2025}
``Reassurance Robots: OCD in the Age of Generative AI,'' arXiv:2602.19401, 2025.

\bibitem{rimal2026personality}
B.~Rimal, ``Can Personality Prompting Change How LLMs Think?'' 2026. [Online]. Available: \url{https://bipinrimal.com.np/blog/013-personality-prompting}

\bibitem{salkovskis1985}
P.~M.~Salkovskis, ``Obsessional-compulsive problems: A cognitive-behavioural analysis,'' \textit{Behaviour Research and Therapy}, vol.~23, no.~5, pp.~571--583, 1985.

\bibitem{serapio2025}
G.~Serapio-Garc\'{i}a, M.~Safdari~\textit{et~al.}, ``A psychometric framework for evaluating and shaping personality traits in large language models,'' \textit{Nature Machine Intelligence}, 2025.

\bibitem{steele1995}
C.~M.~Steele and J.~Aronson, ``Stereotype threat and the intellectual test performance of African Americans,'' \textit{J.\ Personality and Social Psychology}, vol.~69, no.~5, pp.~797--811, 1995.

\bibitem{twenge2018}
J.~M.~Twenge, T.~E.~Joiner, M.~L.~Rogers, and G.~N.~Martin, ``Increases in Depressive Symptoms, Suicide-Related Outcomes, and Suicide Rates Among U.S.\ Adolescents After 2010,'' \textit{Clinical Psychological Science}, vol.~6, no.~1, pp.~3--17, 2018.

\bibitem{aitm2025}
``Red-Teaming LLM Multi-Agent Systems via Communication Attacks,'' arXiv:2502.14847, in \textit{Findings of ACL}, 2025.

\bibitem{darkside2025}
``The Dark Side of LLMs: Agent-based Attacks for Complete Computer Takeover,'' arXiv:2507.06850, 2025.

\bibitem{bullying2025}
``Bullying the Machine: How Personas Increase LLM Vulnerability,'' arXiv:2505.12692, 2025.

\bibitem{psychogenic2025}
J.~Au~Yeung~\textit{et~al.}, ``The Psychogenic Machine: Simulating AI Psychosis, Delusion Reinforcement and Harm Enablement in Large Language Models,'' arXiv:2509.10970, 2025.

\bibitem{nato2025}
NATO Allied Command Transformation, ``Cognitive Warfare 2025: Chief Scientist Report,'' 2025.

\bibitem{shen2024}
X.~Shen~\textit{et~al.}, ```Do Anything Now': Characterizing and Evaluating In-The-Wild Jailbreak Prompts on Large Language Models,'' in \textit{Proc.\ ACM CCS}, 2024.

\end{thebibliography}

\end{document}
